\section{Puntos de acumulación}

\begin{definition}[Punto de acumulación]
  Sea \(X \subset \R^n\). Un punto \(x \in \R^n\) (no necesariamente de \(X\)) es un \emph{punto de acumulación} de \(X\) si \(x \in \cl{X \setminus \{x\}}\).
\end{definition}

\begin{definition}[Conjunto derivado]
  El \emph{conjunto derivado} de \(X\) es el conjunto \(X'\) de todos los puntos de acumulación de \(X\).
\end{definition}

\begin{prop}
  Dados \(X \subset \R^n\) y \(x \in \R^n\), son equivalentes:
  \begin{enumerate}[label=(\arabic*)]
    \item \(x\) es un punto de acumulación de \(X\).
    \item Para todo \(r > 0\), \(B_r(x)\) contiene un punto de \(X\) distinto de \(x\).
    \item Todo entorno abierto de \(x\) contiene un punto de \(X\) distinto de \(x\).
    \item Todo entorno de \(x\) contiene un punto de \(X\) distinto de \(x\).
  \end{enumerate}
\end{prop}

\begin{proof}
  Es la caracterización de la clausura (proposición de la sección de clausura) aplicada al conjunto \(X \setminus \{x\}\).
\end{proof}

\begin{prop}
  Si \(X \subset \R^n\), entonces \(\cl{X} = X \cup X'\).
\end{prop}

\begin{proof}
  (\(\supset\)) Se tiene \(X \subset \cl{X}\). Si \(x \in X'\), entonces \(x \in \cl{X \setminus \{x\}} \subset \cl{X}\).

  (\(\subset\)) Sea \(x \in \cl{X}\) y supongamos \(x \notin X\). Entonces \(X \setminus \{x\} = X\), de modo que \(x \in \cl{X} = \cl{X \setminus \{x\}}\), es decir, \(x \in X'\).
\end{proof}

\begin{prop}
  Sea \(X \subset \R^n\). Un punto \(x\) es de acumulación de \(X\) si y sólo si todo entorno abierto de \(x\) contiene infinitos puntos de \(X\).
\end{prop}

\begin{proof}
  (\(\Leftarrow\)) Si todo entorno abierto de \(x\) contiene infinitos puntos de \(X\), contiene alguno distinto de \(x\), y entonces \(x \in X'\).

  (\(\Rightarrow\)) Sea \(x \in X'\) y \(N\) un entorno abierto de \(x\); existe \(r > 0\) con \(B_r(x) \subset N\). Construimos recursivamente puntos \(x_1, x_2, \dots \in N \cap X\), todos distintos de \(x\), con \(d(x_1, x) > d(x_2, x) > \cdots > 0\). Como \(x \in X'\), \(B_r(x)\) contiene algún \(x_1 \in X\) con \(x_1 \neq x\). Dado \(x_m\), la bola \(B_{d(x_m, x)}(x)\) es un entorno abierto de \(x\), y contiene algún \(x_{m+1} \in X\) distinto de \(x\); entonces \(0 < d(x_{m+1}, x) < d(x_m, x)\). Como las distancias a \(x\) son estrictamente decrecientes, los \(x_m\) son distintos dos a dos. Luego \(N \cap X\) es infinito.
\end{proof}

\begin{prop}
  Si \(X \subset \R^n\), entonces \({(X')}' \subset X'\).
\end{prop}

\begin{proof}
  Sea \(x \in {(X')}'\) y \(N\) un entorno abierto de \(x\). Existe \(y \in X'\) con \(y \in N\) e \(y \neq x\). Como \(N\) es un entorno abierto de \(y\) y \(y \in X'\), \(N\) contiene infinitos puntos de \(X\). Por la proposición anterior, \(x \in X'\).
\end{proof}

\begin{corollary}
  El conjunto derivado de todo subconjunto de \(\R^n\) es cerrado.
\end{corollary}

\begin{proof}
  \(\cl{X'} = X' \cup {(X')}' = X'\), luego \(X'\) es cerrado.
\end{proof}

\begin{definition}[Conjunto perfecto]
  Un conjunto \(X \subset \R^n\) es \emph{perfecto} si \(X = X'\).
\end{definition}

\begin{definition}[Punto aislado]
  Sea \(X \subset \R^n\). Un punto \(x \in X\) es \emph{aislado} en \(X\) si existe un entorno \(N\) de \(x\) tal que \(N \cap X = \{x\}\).
\end{definition}

\begin{definition}[Conjunto discreto]
  Un conjunto es \emph{discreto} si todos sus puntos son aislados.
\end{definition}

\begin{eg}
  \(\Z \subset \R\) es discreto: para \(k \in \Z\), el entorno \(\left( k - \frac{1}{2}, k + \frac{1}{2} \right)\) cumple \(\left( k - \frac{1}{2}, k + \frac{1}{2} \right) \cap \Z = \{k\}\).
\end{eg}

\section{Sucesiones en varias dimensiones}

En esta sección trabajamos en \(\R^m\), para no confundir la dimensión con el índice de las sucesiones. Una sucesión \({(x_n)}_{n \in \N} \subset \R^m\) \emph{converge} a \(L \in \R^m\) si
\[
  (\forall \e > 0)(\exists n_0 \in \N) \colon (\forall n > n_0) \quad d(x_n, L) < \e.
\]

\begin{lemma}
  Si una sucesión en \(\R^m\) tiene límite, éste es único.
\end{lemma}

\begin{proof}
  Si \(x_n \to L\) y \(x_n \to L'\), dado \(\e > 0\) existe \(n_0\) tal que para \(n > n_0\) vale \(d(x_n, L) < \e\) y \(d(x_n, L') < \e\). Entonces \(d(L, L') \leq d(L, x_n) + d(x_n, L') < 2\e\). Como \(\e\) es arbitrario, \(d(L, L') = 0\) y \(L = L'\).
\end{proof}

\begin{prop}
  Sean \({(x_n)}_{n \in \N} \subset \R^m\) y \(L \in \R^m\). Son equivalentes:
  \begin{enumerate}[label=(\arabic*)]
    \item \(L = \lim_{n \to \infty} x_n\).
    \item Para todo \(r > 0\) existe \(n_0\) tal que \(x_n \in B_r(L)\) para todo \(n > n_0\).
    \item Para todo entorno abierto \(N\) de \(L\) existe \(n_0\) tal que \(x_n \in N\) para todo \(n > n_0\).
    \item Para todo entorno \(N\) de \(L\) existe \(n_0\) tal que \(x_n \in N\) para todo \(n > n_0\).
  \end{enumerate}
\end{prop}

\begin{proof}
  (4) \(\Rightarrow\) (3) \(\Rightarrow\) (2) \(\Rightarrow\) (1) son inmediatos de las definiciones. (1) \(\Rightarrow\) (4): sea \(N\) un entorno de \(L\). Como \(L \in N^{\circ}\), existe \(r > 0\) con \(B_r(L) \subset N^{\circ} \subset N\). Por (1) existe \(n_0\) con \(d(L, x_n) < r\) para \(n > n_0\), es decir, \(x_n \in B_r(L) \subset N\).
\end{proof}

\begin{prop}
  Toda sucesión convergente en \(\R^m\) es acotada.
\end{prop}

\begin{proof}
  Sea \(x_n \to x\). Existe \(n_0\) tal que \(d(x_n, x) < 1\) para \(n > n_0\), y entonces \(\norm{x_n} \leq \norm{x_n - x} + \norm{x} < 1 + \norm{x}\). Luego \(\norm{x_n} \leq M\) para todo \(n\), donde \(M = \max\{ \norm{x_1}, \dots, \norm{x_{n_0}}, 1 + \norm{x} \}\).
\end{proof}

\begin{prop}
  Si \(a_n \to a\) y \(b_n \to b\) en \(\R^m\), entonces \(a_n + b_n \to a + b\).
\end{prop}

\begin{proof}
  Dado \(\e > 0\), existen \(n_0\) y \(n_1\) tales que \(d(a_n, a) < \frac{\e}{2}\) para \(n > n_0\) y \(d(b_n, b) < \frac{\e}{2}\) para \(n > n_1\). Para \(n > \max\{n_0, n_1\}\),
  \[
    d(a_n + b_n, a + b) = \norm{(a_n - a) + (b_n - b)} \leq d(a_n, a) + d(b_n, b) < \e. \qedhere
  \]
\end{proof}

\begin{prop}
  Si \(x_n \to x\) en \(\R^m\) y \(\lambda_n \to \lambda\) en \(\R\), entonces \(\lambda_n x_n \to \lambda x\).
\end{prop}

\begin{proof}
  Como ambas sucesiones convergen, son acotadas: existe \(M > 0\) con \(\norm{x_n} \leq M\) para todo \(n\) y \(|\lambda| \leq M\). Dado \(\e > 0\), existen \(n_0\) y \(n_1\) tales que \(\norm{x_n - x} < \frac{\e}{2M}\) para \(n > n_0\) y \(|\lambda_n - \lambda| < \frac{\e}{2M}\) para \(n > n_1\). Para \(n > \max\{n_0, n_1\}\),
  \[
    \norm{\lambda_n x_n - \lambda x} \leq \norm{x_n} \, |\lambda_n - \lambda| + |\lambda| \, \norm{x_n - x} < M \cdot \frac{\e}{2M} + M \cdot \frac{\e}{2M} = \e. \qedhere
  \]
\end{proof}

\begin{prop}
  Sea \({(x_n)}_{n \in \N} \subset \R^m\), con \(x_n = \bigl( x_n^{(1)}, \dots, x_n^{(m)} \bigr)\). Son equivalentes:
  \begin{enumerate}[label=(\arabic*)]
    \item La sucesión \({(x_n)}\) converge en \(\R^m\).
    \item Para todo \(i \in \{1, \dots, m\}\), la sucesión \({\bigl( x_n^{(i)} \bigr)}_{n \in \N}\) converge en \(\R\).
  \end{enumerate}
  Además, en ese caso, \(\lim_{n \to \infty} x_n = (L_1, \dots, L_m)\), donde \(L_i = \lim_{n \to \infty} x_n^{(i)}\).
\end{prop}

\begin{proof}
  (1) \(\Rightarrow\) (2) Sea \(x_n \to L = (L_1, \dots, L_m)\) y \(\e > 0\). Existe \(n_0\) tal que \(\norm{x_n - L} < \e\) para \(n > n_0\). Para cada \(i\),
  \[
    \bigl| x_n^{(i)} - L_i \bigr| \leq \sqrt{\sum_{j=1}^{m} {\bigl( x_n^{(j)} - L_j \bigr)}^2} = \norm{x_n - L} < \e,
  \]
  de modo que \(x_n^{(i)} \to L_i\).

  (2) \(\Rightarrow\) (1) Sea \(x_n^{(i)} \to L_i\) para cada \(i\), y \(L = (L_1, \dots, L_m)\). Dado \(\e > 0\), para cada \(i\) existe \(n_i\) con \(\bigl| x_n^{(i)} - L_i \bigr| < \frac{\e}{\sqrt{m}}\) para \(n > n_i\). Sea \(n_0 = \max\{n_1, \dots, n_m\}\). Para \(n > n_0\),
  \[
    \norm{x_n - L} = \sqrt{\sum_{i=1}^{m} {\bigl( x_n^{(i)} - L_i \bigr)}^2} < \sqrt{m \cdot \frac{\e^2}{m}} = \e. \qedhere
  \]
\end{proof}

\section{Conjuntos compactos}

\begin{theorem}
  Un subconjunto \(F \subset \R^m\) es cerrado si y sólo si para toda sucesión \({(x_n)} \subset F\) tal que \(x_n \to x\) se cumple \(x \in F\).
\end{theorem}

\begin{proof}
  (\(\Rightarrow\)) Sea \(F\) cerrado, \({(x_n)} \subset F\) y \(x_n \to x\). Si fuera \(x \notin F\), entonces \(x \in \R^m \setminus F\), que es abierto, y existiría \(r > 0\) con \(B_r(x) \subset \R^m \setminus F\). Como \(x_n \to x\), existe \(n_0\) con \(x_n \in B_r(x)\) para \(n > n_0\), y entonces \(x_n \notin F\), lo cual es absurdo. Luego \(x \in F\).

  (\(\Leftarrow\)) Supongamos que \(F\) no es cerrado. Entonces \(\cl{F} \neq F\), y existe \(x \in \cl{F} \setminus F\). Para cada \(n \in \N\), por la caracterización de la clausura existe \(x_n \in F \cap B_{1/n}(x)\). Entonces \({(x_n)} \subset F\) y \(x_n \to x \notin F\), lo que contradice la hipótesis. Luego \(F\) es cerrado.
\end{proof}

\begin{eg}
  El intervalo \((0, 1]\) no es cerrado: \(\frac{1}{n} \in (0, 1]\) para todo \(n\) y \(\frac{1}{n} \to 0 \notin (0, 1]\).
\end{eg}

\begin{definition}[Relativamente compacto]
  Un subconjunto \(F \subset \R^m\) es \emph{relativamente compacto} si toda sucesión en \(F\) tiene una subsucesión convergente (en \(\R^m\)).
\end{definition}

\begin{definition}[Compacto]
  Un subconjunto \(F \subset \R^m\) es \emph{compacto} si toda sucesión en \(F\) tiene una subsucesión que converge a un punto de \(F\). Equivalentemente, \(F\) es relativamente compacto y cerrado.
\end{definition}

\begin{note}
  La equivalencia se debe a que, si \(F\) es compacto, toda sucesión convergente en \(F\) converge a un punto de \(F\) (pues una subsucesión lo hace), de modo que \(F\) es cerrado por el teorema anterior. Recíprocamente, si \(F\) es cerrado y relativamente compacto, el límite de la subsucesión convergente pertenece a \(F\) por el mismo teorema.
\end{note}

\begin{prop}
  Un subconjunto \(F \subset \R^m\) es relativamente compacto si y sólo si es acotado.
\end{prop}

\begin{proof}
  (\(\Rightarrow\)) Por contrarrecíproco, supongamos que \(F\) no es acotado. Para cada \(n \in \N\) existe \(y_n \in F\) con \(\norm{y_n} > n\). Si \(f\colon \N \to \N\) es estrictamente creciente, \(\norm{y_{f(n)}} > f(n) \geq n\), de modo que \({(y_{f(n)})}\) no es acotada y por lo tanto no converge. Luego ninguna subsucesión de \({(y_n)}\) converge, y \(F\) no es relativamente compacto.

  (\(\Leftarrow\)) Sea \(F\) acotado y \({(x_k)}_{k \in \N} \subset F\). Como \(\norm{x_k} \leq K\) para algún \(K\), cada sucesión de coordenadas \({\bigl( x_k^{(i)} \bigr)}_{k}\) es acotada en \(\R\). Por el teorema de Bolzano--Weierstrass en \(\R\) existe una función estrictamente creciente \(h_1\colon \N \to \N\) tal que \({\bigl( x_{h_1(k)}^{(1)} \bigr)}_{k}\) converge. Del mismo modo, como \({\bigl( x_{h_1(k)}^{(2)} \bigr)}_{k}\) es acotada, existe \(g_2\) estrictamente creciente tal que, con \(h_2 = h_1 \circ g_2\), la sucesión \({\bigl( x_{h_2(k)}^{(2)} \bigr)}_{k}\) converge; la primera coordenada de \({(x_{h_2(k)})}\) sigue convergiendo por ser subsucesión de una convergente. Repitiendo \(m\) veces se obtiene una subsucesión \({(x_{h_m(k)})}_{k}\) en la cual todas las coordenadas convergen, y por la proposición anterior converge en \(\R^m\).
\end{proof}

\begin{corollary}[Heine--Borel]
  Un subconjunto \(F \subset \R^m\) es compacto si y sólo si es cerrado y acotado.
\end{corollary}

\begin{proof}
  Es consecuencia de la definición de compacto y de la proposición anterior.
\end{proof}

\section{Cubrimientos y compacidad}

Los resultados siguientes corresponden a la clase siguiente, pero se agrupan aquí para reunir lo referente a conjuntos compactos. En esta sección volvemos a trabajar en \(\R^n\).

\begin{definition}[Cubrimiento]
  Sea \(X \subset \R^n\). Una familia \(\mathscr{U}\) de subconjuntos de \(\R^n\) es un \emph{cubrimiento} de \(X\) si \(X \subset \bigcup_{U \in \mathscr{U}} U\). Es un \emph{cubrimiento abierto} si todos sus elementos son abiertos.
\end{definition}

\begin{definition}[Subcubrimiento]
  Un \emph{subcubrimiento} de \(\mathscr{U}\) es una subfamilia \(\mathscr{V} \subset \mathscr{U}\) que también es un cubrimiento de \(X\).
\end{definition}

\begin{theorem}[Lindelöf]
  Todo cubrimiento abierto de un subconjunto de \(\R^n\) admite un subcubrimiento a lo sumo numerable.
\end{theorem}

\begin{proof}
  Sea \(X \subset \R^n\) y \(\mathscr{U}\) un cubrimiento abierto de \(X\). Si \(X = \varnothing\), la familia vacía es un subcubrimiento. Si \(X \neq \varnothing\), entonces \(\mathscr{U} \neq \varnothing\) y fijamos \(U_0 \in \mathscr{U}\).

  Para cada \(q \in \Q^n\) y cada \(s \in \Q\), \(s > 0\), elegimos \(U(q, s) \in \mathscr{U}\) así: si existen elementos de \(\mathscr{U}\) que contienen a \(B_s(q)\), elegimos uno cualquiera; si no, ponemos \(U(q, s) = U_0\).

  Sea \(x \in X\). Existen \(U \in \mathscr{U}\) con \(x \in U\) y \(r > 0\) con \(B_r(x) \subset U\). Como \(\Q^n\) es denso en \(\R^n\), existe \(q \in \Q^n\) con \(d(x, q) < \frac{r}{2}\), y existe \(s \in \Q\) con \(d(x, q) < s < \frac{r}{2}\). Entonces \(x \in B_s(q)\), y si \(y \in B_s(q)\), \(d(y, x) \leq d(y, q) + d(q, x) < 2s < r\), es decir, \(B_s(q) \subset B_r(x) \subset U\). Luego hay elementos de \(\mathscr{U}\) que contienen a \(B_s(q)\), de modo que \(x \in B_s(q) \subset U(q, s)\).

  Esto prueba que \(\mathscr{V} = \{ U(q, s) : q \in \Q^n, \ s \in \Q, \ s > 0 \}\) es un subcubrimiento de \(\mathscr{U}\). Es a lo sumo numerable porque es la imagen de \(\Q^n \times (\Q \cap (0, \infty))\), que es numerable, por la función \((q, s) \mapsto U(q, s)\).
\end{proof}

\begin{theorem}[Heine--Borel--Lebesgue]
  Un subconjunto \(F \subset \R^n\) es compacto si y sólo si todo cubrimiento abierto de \(F\) admite un subcubrimiento finito.
\end{theorem}

\begin{proof}
  (\(\Rightarrow\)) Sea \(F\) compacto y \(\mathscr{U}\) un cubrimiento abierto de \(F\). Por el teorema de Lindelöf existe un subcubrimiento \(\mathscr{V}\) a lo sumo numerable; si es finito no hay nada que probar, de modo que supongamos \(\mathscr{V} = \{ V_1, V_2, \dots \}\) numerable. Supongamos, por el absurdo, que \(\mathscr{U}\) no admite subcubrimiento finito. En particular, para cada \(m\) la familia \(\{ V_1, \dots, V_m \}\) no cubre \(F\), y existe \(x_m \in F \setminus \bigcup_{i=1}^{m} V_i\). Como \(F\) es compacto, existe una subsucesión \(x_{h(m)} \to L \in F\) con \(h\) estrictamente creciente. Como \(\mathscr{V}\) cubre \(F\), existe \(s\) con \(L \in V_s\). Como \(V_s\) es abierto, existe \(m_0\) tal que \(x_{h(m)} \in V_s\) para \(m > m_0\). Pero para \(m > \max\{m_0, s\}\) se tiene \(h(m) \geq m > s\), y por construcción \(x_{h(m)} \notin V_i\) para todo \(i \leq h(m)\), en particular \(x_{h(m)} \notin V_s\). Absurdo.

  (\(\Leftarrow\)) Supongamos que todo cubrimiento abierto de \(F\) admite un subcubrimiento finito, pero que \(F\) no es compacto. Entonces existe \({(x_m)} \subset F\) sin subsucesiones convergentes a puntos de \(F\). Sea \(x \in F\). Como \(x\) no es límite de ninguna subsucesión, por el mismo argumento que en \(\R\) (un punto es límite de una subsucesión si y sólo si todo entorno contiene infinitos términos) existe \(r_x > 0\) tal que \(S_x = \{ m \in \N : x_m \in B_{r_x}(x) \}\) es finito. La familia \(\{ B_{r_x}(x) : x \in F \}\) es un cubrimiento abierto de \(F\), y por hipótesis existen \(y_1, \dots, y_k \in F\) con \(F \subset \bigcup_{i=1}^{k} B_{r_{y_i}}(y_i)\). Como \(x_m \in F\) para todo \(m\), cada \(m \in \N\) pertenece a algún \(S_{y_i}\), es decir, \(\N \subset S_{y_1} \cup \cdots \cup S_{y_k}\), que es finito. Absurdo.
\end{proof}

\begin{prop}[Número de Lebesgue]
  Sea \(K \subset \R^n\) compacto y \(\mathscr{U}\) un cubrimiento abierto de \(K\). Entonces existe \(\delta > 0\) tal que si \(x, y \in K\) y \(d(x, y) < \delta\), existe \(U \in \mathscr{U}\) con \(x, y \in U\). Un tal \(\delta\) se llama \emph{número de Lebesgue} del cubrimiento.
\end{prop}

\begin{proof}
  Para cada \(x \in K\) elegimos \(U_x \in \mathscr{U}\) con \(x \in U_x\) y \(r_x > 0\) con \(B_{r_x}(x) \subset U_x\). La familia \(\bigl\{ B_{r_x/2}(x) : x \in K \bigr\}\) es un cubrimiento abierto de \(K\), y por el teorema anterior existen \(x_1, \dots, x_k \in K\) con \(K \subset \bigcup_{i=1}^{k} B_{r_{x_i}/2}(x_i)\). Sea \(\delta = \frac{1}{2} \min\{ r_{x_1}, \dots, r_{x_k} \}\).

  Sean \(x, y \in K\) con \(d(x, y) < \delta\). Existe \(i\) con \(x \in B_{r_{x_i}/2}(x_i)\), y entonces
  \[
    d(x_i, y) \leq d(x_i, x) + d(x, y) < \frac{r_{x_i}}{2} + \delta \leq \frac{r_{x_i}}{2} + \frac{r_{x_i}}{2} = r_{x_i}.
  \]
  Luego \(x, y \in B_{r_{x_i}}(x_i) \subset U_{x_i} \in \mathscr{U}\).
\end{proof}
